% GENERATED FILE. Edit canonical lessons, then run npm run generate:books.
% canonical-source-hash: fnv1a64:f0e30cc9050aa2da
% canonical-lessons: UR-C43-ammi, UR-C43-beta, UR-C43-beti, UR-C43-dada, UR-C43-dadi

\chapter{Mum, Son, Daughter, Grandfather, Grandmother}
\label{ch:ur-mum-son-daughter-grandfather-grandmother}
\hlchaptermodality{43}

\begin{chapteropening}
\textbf{By the end of this chapter:} \emph{I can say mum, a son, a daughter, a grandfather and a grandmother.}

Complete the last lesson of chapter 43: I can say mum, a son, a daughter, a grandfather and a grandmother.
\end{chapteropening}

\section[\texorpdfstring{\ur{امی} (ammī)}{امی (ammī)}]{\ur{امی} (ammī) — mum}

\label{lesson:UR-C43-ammi}



\begin{quote}
Before the new one: say the Urdu for a father, then the Urdu for dad.
\end{quote}

\subsection*{You'll want to know: \ur{امی}}
\textbf{\ur{امی}} — \emph{ammī} — \textquotedblleft{}mum\textquotedblright{}.

\begin{cousinweb}
\textbf{\ur{امی}} (\emph{ammī}) is the fond form of \emph{ammā}, \textquotedblleft{}mother\textquotedblright{}, which CDIAL lists under Sanskrit \emph{ambā}, \textquotedblleft{}mother\textquotedblright{} (CDIAL 574). CDIAL lists \emph{ammā}, not \emph{ammī}.
\end{cousinweb}

\begin{grammarlens}[title={What you've built}]
One more word for this chapter.
\end{grammarlens}

\subsection*{Your turn}
\begin{itemize}
\raggedright
  \item \emph{Say it:} \emph{ammī}
  \item \emph{Say it:} \emph{ammī}, once more
  \item \emph{Say it:} say \emph{ammī}
  \item \emph{Recall:} say the Urdu for a father, then the Urdu for dad, then say \emph{ammī} again
\end{itemize}

\begin{tcolorbox}[breakable,colback=teal!4,colframe=teal!35!black,title={Before you move on}]
What does \ur{امی} mean? (\textquotedblleft{}Mum\textquotedblright{}.) Say it once more.
\end{tcolorbox}
\section[\texorpdfstring{\ur{بیٹا} (beṭā)}{بیٹا (beṭā)}]{\ur{بیٹا} (beṭā) — a son}

\label{lesson:UR-C43-beta}



\begin{quote}
Before the new one: say the Urdu for dad, then the Urdu for mum.
\end{quote}

\subsection*{You'll want to know: \ur{بیٹا}}
\textbf{\ur{بیٹا}} — \emph{beṭā} — \textquotedblleft{}a son\textquotedblright{}.

\begin{cousinweb}
\textbf{\ur{بیٹا}} (\emph{beṭā}) is a tadbhava, a word inherited and reshaped by sound change. It comes from \emph{*beṭṭa}, \textquotedblleft{}boy\textquotedblright{}, a form CDIAL reconstructs from its descendants (CDIAL 9238). CDIAL lists it in an entry headed \emph{*biḍḍa}, \textquotedblleft{}defective\textquotedblright{}, beside Prakrit \emph{biṭṭa}, \textquotedblleft{}boy\textquotedblright{}; no Sanskrit source is agreed.
\end{cousinweb}

\begin{grammarlens}[title={What you've built}]
One more word for this chapter.
\end{grammarlens}

\subsection*{Your turn}
\begin{itemize}
\raggedright
  \item \emph{Say it:} \emph{beṭā}
  \item \emph{Say it:} \emph{beṭā}, once more
  \item \emph{Say it:} say \emph{beṭā}
  \item \emph{Recall:} say the Urdu for dad, then the Urdu for mum, then say \emph{beṭā} again
\end{itemize}

\begin{tcolorbox}[breakable,colback=teal!4,colframe=teal!35!black,title={Before you move on}]
What does \ur{بیٹا} mean? (\textquotedblleft{}A son\textquotedblright{}.) Say it once more.
\end{tcolorbox}
\section[\texorpdfstring{\ur{بیٹی} (beṭī)}{بیٹی (beṭī)}]{\ur{بیٹی} (beṭī) — a daughter}

\label{lesson:UR-C43-beti}



\begin{quote}
Before the new one: say the Urdu for mum, then the Urdu for a son.
\end{quote}

\subsection*{You'll want to know: \ur{بیٹی}}
\textbf{\ur{بیٹی}} — \emph{beṭī} — \textquotedblleft{}a daughter\textquotedblright{}.

\begin{cousinweb}
\textbf{\ur{بیٹی}} (\emph{beṭī}) is the feminine of \emph{beṭā}, \textquotedblleft{}son\textquotedblright{}. CDIAL lists both under a reconstructed \emph{*beṭṭa}, in an entry headed \emph{*biḍḍa}, \textquotedblleft{}defective\textquotedblright{}, beside Prakrit \emph{biṭṭa}, \textquotedblleft{}boy\textquotedblright{} (CDIAL 9238). No Sanskrit source is agreed.
\end{cousinweb}

\begin{grammarlens}[title={What you've built}]
One more word for this chapter.
\end{grammarlens}

\subsection*{Your turn}
\begin{itemize}
\raggedright
  \item \emph{Say it:} \emph{beṭī}
  \item \emph{Say it:} \emph{beṭī}, once more
  \item \emph{Say it:} say \emph{beṭī}
  \item \emph{Recall:} say the Urdu for mum, then the Urdu for a son, then say \emph{beṭī} again
\end{itemize}

\begin{tcolorbox}[breakable,colback=teal!4,colframe=teal!35!black,title={Before you move on}]
What does \ur{بیٹی} mean? (\textquotedblleft{}A daughter\textquotedblright{}.) Say it once more.
\end{tcolorbox}
\section[\texorpdfstring{\ur{دادا} (dādā)}{دادا (dādā)}]{\ur{دادا} (dādā) — a grandfather (father's father)}

\label{lesson:UR-C43-dada}



\begin{quote}
Before the new one: say the Urdu for a son, then the Urdu for a daughter.
\end{quote}

\subsection*{You'll want to know: \ur{دادا}}
\textbf{\ur{دادا}} — \emph{dādā} — \textquotedblleft{}a grandfather (father's father)\textquotedblright{}.

\begin{cousinweb}
\textbf{\ur{دادا}} (\emph{dādā}) is a tadbhava, a word inherited and reshaped by sound change. It comes from \emph{*dādda}, \textquotedblleft{}father or other elderly relative\textquotedblright{}, a form CDIAL reconstructs from its descendants (CDIAL 6261).
\end{cousinweb}

\begin{grammarlens}[title={What you've built}]
One more word for this chapter.
\end{grammarlens}

\subsection*{Your turn}
\begin{itemize}
\raggedright
  \item \emph{Say it:} \emph{dādā}
  \item \emph{Say it:} \emph{dādā}, once more
  \item \emph{Say it:} say \emph{dādā}
  \item \emph{Recall:} say the Urdu for a son, then the Urdu for a daughter, then say \emph{dādā} again
\end{itemize}

\begin{tcolorbox}[breakable,colback=teal!4,colframe=teal!35!black,title={Before you move on}]
What does \ur{دادا} mean? (\textquotedblleft{}A grandfather (father's father)\textquotedblright{}.) Say it once more.
\end{tcolorbox}
\section[\texorpdfstring{\ur{دادی} (dādī)}{دادی (dādī)}]{\ur{دادی} (dādī) — a grandmother (father's mother)}

\label{lesson:UR-C43-dadi}



\begin{quote}
Before the new one: say the Urdu for a daughter, then the Urdu for a grandfather.
\end{quote}

\subsection*{You'll want to know: \ur{دادی}}
\textbf{\ur{دادی}} — \emph{dādī} — \textquotedblleft{}a grandmother (father's mother)\textquotedblright{}.

\begin{cousinweb}
\textbf{\ur{دادی}} (\emph{dādī}) is a tadbhava, a word inherited and reshaped by sound change. It comes from \emph{*dādda}, \textquotedblleft{}father or other elderly relative\textquotedblright{}, a form CDIAL reconstructs from its descendants (CDIAL 6261). \emph{dādī} is the feminine of \emph{dādā}, \textquotedblleft{}grandfather\textquotedblright{}.
\end{cousinweb}

\begin{grammarlens}[title={What you've built}]
One more word for this chapter.
\end{grammarlens}

\subsection*{Your turn}
\begin{itemize}
\raggedright
  \item \emph{Say it:} \emph{dādī}
  \item \emph{Say it:} \emph{dādī}, once more
  \item \emph{Say it:} say \emph{dādī}
  \item \emph{Recall:} say the Urdu for a daughter, then the Urdu for a grandfather, then say \emph{dādī} again
\end{itemize}

\begin{tcolorbox}[breakable,colback=teal!4,colframe=teal!35!black,title={Before you move on}]
What does \ur{دادی} mean? (\textquotedblleft{}A grandmother (father's mother)\textquotedblright{}.) Say it once more.
\end{tcolorbox}
